\section{A finite-character decomposition}
\label{ub:section}

We construct a decomposition that attains the proposed upper bound.
Let $k\geq 1$, let $a_i\geq b_i\geq 1$ be integers, and set
\[
 F=\prod_{i=1}^k x_i^{a_i}y_i^{b_i},\qquad
 d_i=a_i+b_i,\qquad q_i=b_i+1,\qquad \delta_i=a_i-b_i.
\]
Write
\begin{equation}
 Q=\prod_{i=1}^k q_i,
 \qquad
 N=Q+\sum_{i=1}^k\delta_i\frac{Q}{q_i}.
 \label{ub:N}
\end{equation}
The construction uses a lattice whose nonzero vectors avoid the box of
possible exponent differences, followed by character orthogonality on
the corresponding finite quotient.

\begin{lemma}
\label{ub:lattice}
Let $\delta=(\delta_1,\ldots,\delta_k)^T$ and
$\mathbf{1}=(1,\ldots,1)^T$. Define
\[
 A=\operatorname{diag}(q_1,\ldots,q_k)+\delta\mathbf{1}^T,
 \qquad \Lambda=A\Z^k.
\]
Then $\det A=N$, so $[\Z^k:\Lambda]=N$, and
\begin{equation}
 \Lambda\cap\prod_{i=1}^k[-a_i,b_i]=\{0\}.
 \label{ub:box-avoidance}
\end{equation}
\end{lemma}

\begin{proof}
The lattice is generated by the \emph{columns} of $A$: its $j$th
generator is $q_j\mathbf{e}_j+\delta$, where
$\mathbf{e}_1,\ldots,\mathbf{e}_k$ are the standard basis vectors.
Expand the determinant multilinearly in these columns. Any term that
chooses $\delta$ from two columns vanishes. The term choosing no such
column is $Q$, and the term choosing it only from column $i$ is
$\delta_i\prod_{j\ne i}q_j$. Hence
\[
 \det A=Q+\sum_{i=1}^k\delta_i\prod_{j\ne i}q_j=N>0.
\]
The index of a nonsingular integer matrix's column lattice is the absolute
value of its determinant: integer row and column operations preserve
both quantities and reduce the matrix to diagonal form, where the claim
is immediate. Thus $[\Z^k:\Lambda]=N$.

Suppose $u=Az$ belongs to the box in \eqref{ub:box-avoidance}, with
$z\in\Z^k$, and put $S=\sum_i z_i$. The column convention gives
\begin{equation}
 u_i=q_i z_i+\delta_i S.
 \label{ub:lattice-coordinates}
\end{equation}
If $S>0$, then $u_i\leq b_i=q_i-1$ and $\delta_i\geq0$ imply
$q_i z_i\leq q_i-1$ for every $i$. Thus every $z_i\leq0$, contrary
to $S>0$.

If $S<0$, write $H=-S\geq1$. Since
$a_i=q_i-1+\delta_i$, the lower bound $u_i\geq-a_i$ gives
\[
 q_i z_i\geq-q_i+1+(H-1)\delta_i\geq-q_i+1.
\]
Integrality now implies $z_i\geq0$ for every $i$, contrary to $S<0$.
Finally, if $S=0$, the upper bounds again imply $z_i\leq0$ for all
$i$, and their zero sum forces $z=0$. This proves
\eqref{ub:box-avoidance}, including the cases in which some
$\delta_i$ vanish.
\end{proof}

\begin{lemma}[Finite character orthogonality]
\label{ub:orthogonality}
Let $H$ be a finite abelian group of order $N$, and let
$\widehat H=\operatorname{Hom}(H,\C^\times)$. Then
$|\widehat H|=N$ and, for every $h\in H$,
\begin{equation}
 \sum_{\chi\in\widehat H}\chi(h)
 =\begin{cases}
     N,&h=0,\\
     0,&h\ne0.
   \end{cases}
 \label{ub:character-sum}
\end{equation}
\end{lemma}

\begin{proof}
Write $H$ as a product of cyclic groups
$\Z/m_1\Z\times\cdots\times\Z/m_s\Z$. A character is determined
by independently choosing an $m_j$th root of unity as the image of
the $j$th cyclic generator. There are therefore
$\prod_j m_j=N$ characters. If $h$ has coordinates $h_j$, the sum
in \eqref{ub:character-sum} factors as
\[
 \prod_{j=1}^s
 \left(\sum_{r=0}^{m_j-1}
       \exp\!\left(\frac{2\pi\mathrm{i}\,r h_j}{m_j}\right)\right).
\]
Each geometric sum equals $m_j$ when $h_j=0$ in $\Z/m_j\Z$ and
equals zero otherwise, proving the formula.
\end{proof}

\begin{proposition}
\label{prop:upper}
The polynomial $F$ admits a decomposition into $N$ nonzero weighted
products of powers of linear forms, with degree $d_i$ in the $i$th
block. In particular, its fixed-block rank is at most $N$.
\end{proposition}

\begin{proof}
Use the column lattice $\Lambda$ of Lemma~\ref{ub:lattice}, and put
\[
 H=\Z^k/\Lambda,\qquad G=\widehat H,
 \qquad \theta_i(\chi)=\chi(\mathbf{e}_i+\Lambda)
 \quad(\chi\in G).
\]
Each $\theta_i(\chi)$ is a nonzero root of unity. We claim the exact
identity
\begin{equation}
 \boxed{
 F=
 \frac{1}{N\prod_{i=1}^k\binom{d_i}{a_i}}
 \sum_{\chi\in G}
 \left(\prod_{i=1}^k\theta_i(\chi)^{-a_i}\right)
 \prod_{i=1}^k\bigl(\theta_i(\chi)x_i+y_i\bigr)^{d_i}.}
 \label{eq:character-decomposition}
\end{equation}
Indeed, let $m=(m_1,\ldots,m_k)$ satisfy $0\leq m_i\leq d_i$.
The coefficient of $\prod_i x_i^{m_i}y_i^{d_i-m_i}$ on the right
of \eqref{eq:character-decomposition} is
\begin{equation}
 \frac{\prod_i\binom{d_i}{m_i}}
      {N\prod_i\binom{d_i}{a_i}}
 \sum_{\chi\in G}\chi(m-a+\Lambda),
 \qquad a=(a_1,\ldots,a_k).
 \label{ub:coefficient}
\end{equation}
Since $d_i-a_i=b_i$, the vector $m-a$ lies in
$\prod_i[-a_i,b_i]$. Lemma~\ref{ub:lattice} says that it belongs
to $\Lambda$ precisely when $m=a$. Character orthogonality therefore
makes \eqref{ub:coefficient} zero for $m\ne a$, while for $m=a$
it equals one. This verifies every coefficient of
\eqref{eq:character-decomposition}.

There are $|G|=N$ summands, and every weight in the displayed formula
is nonzero. The normalized support tuples are also distinct: the images
of the standard basis generate $H$, so two characters with the same
$\theta_i$ for every $i$ coincide. Moreover, the product
$\prod_i(\theta_i x_i+y_i)^{d_i}$ has all-$y$ coefficient one and
coefficient $d_i\theta_i$ at
$x_i y_i^{d_i-1}\prod_{j\ne i}y_j^{d_j}$; thus distinct tuples give
nonproportional products. The identity is consequently an $N$-term
fixed-block decomposition over $\C$, and proves the upper bound.
\end{proof}
